\begin{ExerciseBlock}
\ExerciseHeading{}{Verständnisfragen}{two_particle_systems_checks}
Entscheiden Sie, welche Aussagen wahr oder falsch sind, und begründen Sie Ihre Entscheidung.
\begin{enumerate}
\item Bei der Beschreibung von Stößen ist es nötig die wirkenden Kräfte $F(t)$ genau zu kennen.\TrueFalseAnswer{falsch}
\item Bei elastischen und inelastischen Stößen bleibt der Gesamtimpuls erhalten, wenn der äußere Kraftstoß vernachlässigbar ist.\TrueFalseAnswer{wahr}
\item Bei elastischen und inelastischen Stößen gilt die Erhaltung der Bewegungsenergie.\TrueFalseAnswer{falsch}
\item Die Summe der Geschwindigkeiten nach einem elastischen Stoß entspricht der Summe der Geschwindigkeit davor.\TrueFalseAnswer{falsch}
\item Durch einen elastischen Stoß kann ein Stoßpartner eine höhere Geschwindigkeit erreichen, als der schnellere der beiden Stoßpartner vor dem Stoß.\TrueFalseAnswer{wahr}
\item Nach einem vollkommen inelastischen Stoß, bewegen sich die Stoßpartner mit der selben Geschwindigkeit fort.\TrueFalseAnswer{wahr}
\item Bei einem vollkommen inelastischen Stoß entspricht der Verlust an kinetischer Energie der anfänglichen kinetischen Energie der Relativbewegung.\TrueFalseAnswer{wahr}
\item Bei einem inelastischen Stoß reduziert sich die kinetische Energie der Schwerpunktsbewegung.\TrueFalseAnswer{falsch}
\item Bei einem vollkommen inelastischen Stoß ohne äußeren Kraftstoß entspricht die gemeinsame Geschwindigkeit nach dem Stoß der Schwerpunktsgeschwindigkeit davor.\TrueFalseAnswer{wahr}
\end{enumerate}
\end{ExerciseBlock}
